\documentclass[11pt,reqno]{amsart}
\usepackage[a4paper,margin=29mm]{geometry}
\usepackage{amssymb,mathtools,booktabs,longtable}
\usepackage[colorlinks=true,linkcolor=blue,citecolor=blue,urlcolor=blue]{hyperref}
\newcommand{\Q}{\mathbb Q}\newcommand{\Z}{\mathbb Z}
\newcommand{\F}{\mathbb F}\newcommand{\OO}{\mathcal O}
\newcommand{\ns}{\mathrm{ns}}\newcommand{\GL}{\operatorname{GL}}
\newcommand{\Gal}{\operatorname{Gal}}\newcommand{\End}{\operatorname{End}}
\newcommand{\Jac}{\operatorname{Jac}}\newcommand{\Sel}{\operatorname{Sel}}
\newcommand{\Cl}{\operatorname{Cl}}\newcommand{\Tr}{\operatorname{Tr}}
\newcommand{\rank}{\operatorname{rank}}\newcommand{\Lie}{\operatorname{Lie}}
\newcommand{\disc}{\operatorname{disc}}\newcommand{\im}{\operatorname{im}}
\newcommand{\Xns}[1]{X_{\ns}^{+}(#1)}
\newtheorem{theorem}{Theorem}[section]
\newtheorem{proposition}[theorem]{Proposition}
\newtheorem{lemma}[theorem]{Lemma}
\newtheorem{corollary}[theorem]{Corollary}
\theoremstyle{definition}\newtheorem{definition}[theorem]{Definition}
\theoremstyle{remark}\newtheorem{remark}[theorem]{Remark}
\numberwithin{equation}{section}
\emergencystretch=3em
\title[Rational points on the non-split Cartan curve of level 121]{Rational points on the non-split Cartan modular curve of level 121}
\author{Filip Najman}
\date{November 2, 2026}
\subjclass[2020]{11G18, 14G05, 11F80, 11G05}
\keywords{Modular curves, Galois representations, Selmer groups, isogeny descent}
\begin{document}
\begin{abstract}
We determine the rational points on the non-split Cartan modular curve
\(\Xns{121}\). It has exactly seven rational points, all of which are
CM. This completes the classification of the \(11\)-adic Galois images
of non-CM elliptic curves over \(\Q\). A cusp descent reduces the
problem to seven curves of genus five over a number field. We determine
their relevant points by an isogeny descent and Selmer-group Chabauty.
The proof uses a local formal-division calculation at points with
singular reduction and does not assume the generalized Riemann hypothesis.
\end{abstract}
\maketitle
\section{Introduction}

Let \(C_{\ns}(121)\) be a non-split Cartan subgroup of
\(\GL_2(\Z/121\Z)\), obtained from an unramified quadratic algebra,
and let \(N_{\ns}(121)\) be its normalizer. We write
\(\Xns{121}\) for the corresponding modular curve. It has genus
511 and parametrizes elliptic curves with a non-split Cartan
structure of level 121, modulo the normalizer.

Rouse, Sutherland and Zureick-Brown classified the possible
\(11\)-adic Galois images of non-CM elliptic curves over \(\Q\),
apart from the images contained in the inverse image of
\(N_{\ns}(121)\) \cite[Theorem 1.6]{RSZB}. The remaining case is
equivalent to determining whether \(\Xns{121}\) has a non-CM
rational point. The natural map
\[
 \Xns{121}\longrightarrow\Xns{11}
\]
has degree 121. The target is an elliptic curve of rank one.
Schoof and Tzanakis determined its rational points with integral
\(j\)-invariant \cite{ST}. However, that result does not exclude
points with nonintegral \(j\)-invariant coming from level 121.

In this paper we determine all rational points on \(\Xns{121}\).
The main result is the following theorem.

\begin{theorem}\label{thm:main}
The curve \(\Xns{121}\) has exactly seven rational points. They
are CM points with discriminants
\[
 -3,\ -4,\ -12,\ -16,\ -27,\ -67,\ -163.
\]
There is exactly one point for each discriminant. There are no
rational cusps and no non-CM rational points.
\end{theorem}

For an elliptic curve \(A/\Q\), write
\(\rho_{A,11^\infty}:G_\Q\to\GL_2(\Z_{11})\) for the action on
its Tate module and \(\rho_{A,121}\) for its reduction modulo 121.
We obtain the following consequence.

\begin{corollary}\label{cor:images}
Let \(A/\Q\) be a non-CM elliptic curve. Then
\(\rho_{A,121}(G_\Q)\) is not contained in a conjugate of
\(N_{\ns}(121)\). In particular, the classification of
\(11\)-adic Galois images of non-CM elliptic curves over \(\Q\)
is complete.
\end{corollary}
\begin{proof}
Such a containment would give a non-CM rational point on
\(\Xns{121}\), contrary to Theorem~\ref{thm:main}. The last
assertion follows from \cite[Theorem 1.6]{RSZB}.
\end{proof}

We prove Theorem~\ref{thm:main} by working with the image on
\(\Xns{11}\). A cusp descent and a sieve at six primes leave
seven possibilities for the descent classes, together with a
congruence modulo \(4356\). Each class is represented by a known
CM point. A further two-descent gives seven auxiliary curves of
genus five over a number field of degree 15. We retain the
condition that their points arise from rational points on
\(\Xns{11}\).

We then apply the isogeny version of Selmer-group Chabauty
\cite[Theorem 2.1 and Remark 2.2]{Stoll}. We compute the relevant
Selmer groups after adjoining \(\zeta_{11}\) and compare their
localizations at 11 with images obtained by repeated formal
division. While a point remains sufficiently deep in the formal
group, injectivity of Selmer localization forces a global
division. The final local image is incompatible with the computed
Selmer image, which excludes an additional rational point.

The local calculation is different in the odd and even classes.
In the even classes the known point has smooth reduction. In the
odd classes it reduces to the singular point of the integral
model, and we work over ramified quadratic completions. We give
the formal-group calculation in both cases. The descent requires
neither a basis of the Mordell-Weil groups of the auxiliary
Jacobians nor a height bound for the image on \(\Xns{11}\).

The number-field computations use an unconditional discriminant
bound to certify the required class-group input. GRH is used only
to find candidate generators, which are then certified by exact
membership and local independence. Thus Theorem~\ref{thm:main}
does not depend on GRH.

\subsection{Notation}
For a number field \(k\), a finite set \(S\) of finite primes and
a prime \(p\), put
\[
 k(S,p)=\{a\in k^\times/k^{\times p}:
 v_\mathfrak q(a)\equiv0\pmod p\text{ for }\mathfrak q\notin S\}.
\]
We use additive notation for this \(\F_p\)-vector space.
Valuations on local fields are normalized to have value group
\(\Z\). If \(\lambda\) is an isogeny, \(\mu_v\) denotes its
local Kummer map.

\subsection{The CM points}
For each discriminant in Theorem~\ref{thm:main}, the CM order
has class number one and 11 is inert. Its action on the torsion
gives a rational point on \(\Xns{121}\). The corresponding
\(j\)-invariants and the images on \(\Xns{11}\) are listed in
Table~\ref{tab:121}.

We shall use uniqueness of the rational CM lift from level 11.
Lifts of a fixed Cartan algebra modulo 11 form an affine space
whose translation space has dimension two over \(\F_{11}\).
The canonical CM lift is the origin. The Cartan acts on the
translation space through the norm-one subgroup of
\(\F_{121}^\times\), which has order 12 and acts freely away from
the origin. Modulo elliptic automorphisms, the nonzero orbit
lengths are 12, 6 or 4. CM reciprocity gives this action, so
only the canonical lift is rational. The finite orbit calculation
is included in the ancillary data.

\section{Cusp descent and the auxiliary curves}\label{sec:121}

The natural map \(\Xns{121}\to\Xns{11}\) has degree 121.
We use the model
\begin{equation}\label{eq:E11}
 E=\Xns{11}:\quad y^2+y=x^3-x^2-7x+10,\qquad
 E(\Q)=\Z G,\quad G=(4,5).
\end{equation}
The model and Mordell--Weil generator agree with \cite{ST}.
The seven known points at level 121 map to the following points.
\begin{table}[ht]\centering
\begin{tabular}{rllr}\toprule
\(n\)&\(nG\)&\(D\)&\(j\)\\\midrule
\(-2\)&\((2,-1)\)&\(-27\)&\(-12288000\)\\
\(-1\)&\((4,-6)\)&\(-16\)&\(287496\)\\
\(0\)&\(O\)&\(-12\)&\(54000\)\\
\(1\)&\((4,5)\)&\(-67\)&\(-147197952000\)\\
\(2\)&\((2,0)\)&\(-4\)&\(1728\)\\
\(3\)&\((5/4,7/8)\)&\(-3\)&\(0\)\\
\(4\)&\((-2,3)\)&\(-163\)&\(-262537412640768000\)\\\bottomrule
\end{tabular}
\caption{The images of the rational points on \(\Xns{121}\).}\label{tab:121}
\end{table}

\subsection{The finite sieve}

Put \(K=\Q(z)\), where
\[
 z^5+z^4-4z^3-3z^2+3z+1=0,
 \qquad u=2-z.
\]
The unique prime above 11 is generated by \(u\). The five cusps are
\(r_i s\), with
\[
 (r_i)=(1,4,3,9,5),\qquad
 s=(-z^4-z^3+2z^2+3z+3,-2z^4-4z^3+6z^2+11z+1).
\]
The point \(s\) has order 11. The cusp divisor is a single closed
point of degree five, so \(E\) has no rational cusps.

Suppose a rational point at level 121 maps to \(P=nG\). At any prime
away from 11 at which \(j(P)\) is nonintegral, its valuation is divisible
by 121. Indeed, after a quadratic twist the associated elliptic curve
is a Tate curve. A nontrivial unipotent inertia element modulo 121
cannot belong to the normalizer of an unramified Cartan: this holds
modulo 11, and an element \(1+11N\), with \(N\) nonzero nilpotent
modulo 11, cannot belong to the Cartan congruence kernel either.
Hence the Tate parameter has valuation divisible by 121.
As the polar multiplicity of \(j\) at each cusp of \(E\) is 11, the
intersection multiplicities of \(P\) with the cusps are divisible by 11
away from 11.

For \(m=1,2,3\), let \(a_{m,i}\) be the representative in
\(\{1,\ldots,10\}\) of \(r_i^m\bmod11\), and choose \(f_m\) with
\[
 \operatorname{div}(f_m)=\sum_i a_{m,i}(r_i s)-22(O),
 \qquad \left.f_m(x/y)^{22}\right|_O=1.
\]
The class of \(f_m(P)\) lies in \(K(S,11)\), where
\(S=\{(u)\}\). The class number of \(K\) is one. A basis for this
five-dimensional space is
\[
 u,\quad z^4-3z^2+1,\quad z^2+z-1,\quad z,
 \quad -z^4-z^3+3z^2+3z.
\]
At \(O\) one uses the normalized leading coefficient. 

The sieve starts with all three classes unknown. Local residue-symbol
equations, with cusp residues retained conservatively, give:
\begin{center}\begin{tabular}{rrrrl}\toprule
prime&order of \(\bar G\)&modulus&states&ranks of each class system\\\midrule
353&11&11&7&3\\
23&33&33&10&3 or 5\\
3301&33&33&7&5\\
12101&1089&1089&12&5\\
15731&726&2178&7&5\\
53&12&4356&7&5\\\bottomrule
\end{tabular}\end{center}
Write \(M=4356=4\cdot9\cdot121\). The final states are exactly
\begin{equation}\label{eq:sieve}
 n\equiv n_0\pmod{4356},\qquad -2\le n_0\le4,
 \qquad [f_m(nG)]=[f_m(n_0G)]\quad(m=1,2,3).
\end{equation}
We compute all six local tables from the functions above.
The ancillary program \texttt{six\_prime\_sieve.py} performs the exhaustive affine
linear algebra and congruence lifting.
We compute the classes of the seven known points exactly in the displayed
unit basis and verify that they are the surviving classes.

Set
\[
 \alpha=x(2s)=z^4-z^3-3z^2+3z+2,\qquad
 \delta_{n_0}=x(n_0G-3s)-\alpha.
\]
There is an exact normalized function identity
\[
 \frac{x(P-3s)-\alpha}{x(3s)-\alpha}
   = f_1(P)^7 f_2(P)^{10} f_3(P)^6 f_T(P)^3 h(P)^{11},
\]
where \(h\in K(E)^\times\) and
\[
 \operatorname{div}(f_T)=11(s)-11(O),\qquad
 \left.f_T(x/y)^{11}\right|_O=1.
\]
We verify the identity, including its constant, in the function field
of \(E\). The class of \(f_T(P)\) is an
unramified descent homomorphism on \(E(K)/11E(K)\).
It follows from \eqref{eq:sieve} that
\begin{equation}\label{eq:cusp}
 x(nG-3s)=\alpha+\delta_{n_0}w^{11},\qquad w\in K.
\end{equation}

Let
\[
 L=K(\theta),\quad \theta^3-\theta^2-7\theta+41/4=0,
 \qquad \beta^2=\alpha-\theta,\quad \beta\in L,
 \qquad F=L(\zeta_{11}).
\]
The degrees of \(L,F\) are 15 and 30. The two-descent map \(x-\theta\)
on \(E(L)\) shows that \eqref{eq:cusp} gives a point on
\begin{equation}\label{eq:genus5}
 \eta v^2=\alpha-\theta+\delta_{n_0}w^{11},\qquad
 \eta=\begin{cases}1&n_0\text{ even},\\4-\theta&n_0\text{ odd}.
 \end{cases}
\end{equation}
Indeed, \(-3s\) is twice a point and \(n\equiv n_0\bmod2\).
Put \(a=\delta_{n_0}/\beta^2\) and \(b=a^5\). The ideal of \(a\)
in \(F\) is the square of its unique prime over 2. In particular it
is a unit at 11. Over \(B=F(\sqrt\eta)\), the substitutions
\[
 X=aw,\qquad Y=\sqrt\eta\,a^5v/\beta
\]
give the unit model
\begin{equation}\label{eq:unit}
 C:\quad Y^2=X^{11}+b^2.
\end{equation}
Its known point \(P_0\) has \(X=a\). We write \(J\) for its Jacobian and \(J^\eta\) for its
quadratic twist. In an even case \(\eta=1\). We choose the known
square root \(v_0\in L\) compatibly with the original elliptic point,
and write \(P_0\) for the resulting point with \(w=1\).


\subsection{The odd Selmer groups}

We first consider \(n_0=-1,1,3\) and \(\eta=4-\theta\). Let
\(\tau\) generate \(\operatorname{Gal}(B/F)\). The twist identifies
\(J^\eta(F)\) with the anti-invariant subgroup of \(J(B)\).
The endomorphisms
\[
 \sigma(X,Y)=(\zeta_{11}X,Y),\qquad \lambda=1-\sigma
\]
commute with twisting. The kernel of \(\lambda\) on \(J\) is generated
by \(T=[(0,b)-\infty]\), which has order 11. The action of \(\tau\)
on \(T\) is trivial, so \(J^\eta[\lambda](F)=0\).
The Kummer map is represented over \(B\) by \(Y-b\). Its values
lie in the anti-invariant part of \(B^\times/B^{\times11}\).

The exact discriminants are
\[
 |\operatorname{disc}F|=2^{20}11^{27},\qquad
 |\operatorname{disc}B|=2^{60}11^{56}.
\]
Both fields have root discriminant less than 19. The row \(b=9\) of
Odlyzko's unconditional Table 4 \cite{Od} gives
\[
 |D_M|>53.047^{r_1}20.710^{2r_2}e^{-24.001}.
\]
For any number field of degree \(d\ge330\), this implies
\[
 |D_M|^{1/d}>20.710\,e^{-24.001/330}
 >20.710\left(1-\frac{24.001}{330}\right)>19.
\]
An everywhere unramified cyclic extension of degree 11 of \(F\) or
\(B\) would have degree at least 330 and the same root discriminant.
Thus
\begin{equation}\label{eq:classes}
 \Cl(F)[11]=\Cl(B)[11]=0.
\end{equation}
We verify the discriminants and the displayed rational inequality
exactly. Odlyzko rounds the lower bounds down and the error constant
up, as explained in \cite{OdDescription}.

Let \(S\) consist of the primes over 2 and 11. Each of \(F\) and \(B\)
has four such primes. Their unit ranks are 14 and 29, and both contain
\(\mu_{11}\). It follows from \eqref{eq:classes} that
\[
 \dim F(S,11)=19,\qquad \dim B(S,11)=34.
\]
Norm is onto, since norm composed with restriction is multiplication by
2. Thus the anti-invariant global Kummer space has dimension 15.

We construct this space using candidate generators from the
degree-30 fields
\[
 L(\sqrt{4-\theta}),\qquad L(\sqrt{-11(4-\theta)}).
\]
For each candidate we check the ideal valuations outside \(S\),
map it into \(B\) and divide it by its conjugate. The resulting
element has norm one. We obtain fifteen independent joint local
images. Since the global space has dimension 15, these classes form
a basis. GRH is used to find candidates, but the membership and
independence checks certify the basis without that assumption.

There are three primes \(v\mid11\) of \(F\), all with
\(F_v\simeq\Q_{11}(\zeta_{11})\). At two of them, denoted \(v_1,v_2\),
\(B/F\) is ramified quadratic. At the third it is unramified quadratic.
It is nonsplit at all three, so \(J^\eta[\lambda](F_v)=0\).
The determinant of \(\lambda\) on the tangent space has valuation five
in \(F_v\), and hence
\[
 \dim_{\F_{11}} J^\eta(F_v)/\lambda J^\eta(F_v)=5.
\]
We generate each local image using anti-invariant divisors
\(6(P-\tau P)\), with \(P\in C(B_v)\), until its dimension is five.
Divisors of degree two are needed to obtain the full local image.

Let \(\mathcal S^+\) be the intersection of the certified global
15-dimensional space with these three local images. Omitting the
condition at 2 makes this a Selmer upper bound, which suffices. We compute that
\begin{equation}\label{eq:selmer}
 \dim\mathcal S^+=1\quad(n_0=-1,1,3),\qquad
 \mathcal S^+\longrightarrow
 \prod_{i=1}^2 B_{v_i}^\times/B_{v_i}^{\times11}
 \text{ is injective}.
\end{equation}
The known point has a nonzero class in each of these lines. Thus the
actual Selmer group and \(J^\eta(F)/\lambda J^\eta(F)\) both have
dimension one. Since the rational \(\lambda\)-torsion is zero and
\(\lambda\) has cyclotomic norm 11, this forces rank ten over \(F\).
The nontrivial automorphism of \(F/L\) conjugates the cyclotomic
action, so its fixed subspace has half this dimension. Hence the corresponding
Jacobian has rank exactly five over \(L\).
The rational-point argument below does not use this rank assertion.


\section{Local division}\label{sec:local}

We first describe the formal group of the unit model \eqref{eq:unit}.
At a place \(v\) over 11, put \(k=F_v\) and let \(l\) be either \(k\) in an even case
or a ramified quadratic completion \(B_v\) in an odd case.
Put \(\pi=1-\zeta_{11}\) and \(e=v_l(\pi)\in\{1,2\}\).
Thus \(v_l(11)=10e\). For \(e=1\) use \(\varpi=\pi\).
For \(e=2\) we use
\[
 \varpi=\frac{\sqrt\eta}{\pi^2},\qquad
 v_l(\varpi)=1,\qquad \tau\varpi=-\varpi.
\]
In either case the residue field is \(\F_{11}\).

We use the proper hyperelliptic model of \eqref{eq:unit}. Its special
fibre is geometrically integral, with a single cusp and a smooth point
at infinity. The formal completion at the identity of its relative
Picard group is a smooth formal group of dimension five, since the
obstruction group \(H^2(\mathcal O)\) vanishes on each fibre. We may
construct its formal coordinates from sums of Abel maps at five smooth
sections with distinct residual \(X\)-coordinates. The cotangent lattice
has basis
\[
 \omega_j=\frac{X^{j-1}\,dX}{2Y},\qquad 1\le j\le5.
\]
Write \(H_m\) for formal points whose coordinates lie in
\(\varpi^m\OO_l\).

\begin{lemma}\label{lem:log}
The formal logarithm and exponential have integral coefficients in
these integral coordinates. Moreover,
\[
 \lambda:H_m\xrightarrow{\sim}H_{m+e}\qquad(m\ge1).
\]
In logarithm coordinates dual to the \(\omega_j\), the action of
\(\lambda\) is diagonal with entries \(1-\zeta_{11}^j\).
\end{lemma}
\begin{proof}
The automorphism \(\sigma\) is the identity on the model modulo the
ideal \((\pi)\), including its nonreduced base change to
\(\OO_l/(\pi)\). Hence every coefficient of the formal power series
for \(1-\sigma\) is divisible by \(\pi\). Its linear part divided by
\(\pi\) is invertible, since
\((1-\zeta_{11}^j)/\pi\equiv j\pmod\varpi\).
Dividing the series by \(\pi\) and applying the formal inverse
function theorem proves the asserted isomorphism.

The norm correspondence \(1+\sigma+\cdots+\sigma^{10}\) is zero on
the Jacobian, since the quotient curve has genus zero. In
\(\Z[\sigma]\), multiplication by 11 is therefore a unit times
\(\lambda^{10}\). Every coefficient of its formal series is divisible
by \(\pi^{10}\), and hence by 11.
The logarithm identity
\(\log([11]t)=11\log(t)\) now proves integrality by induction on total
degree. In degree \(d\ge2\), the unknown coefficient is multiplied
by \(11^d-11\), whose quotient by 11 is a unit. The already known
terms are all divisible by 11. The inverse series is integral as well.
Finally \(\sigma^*\omega_j=\zeta_{11}^j\omega_j\), which gives the
diagonal action.
\end{proof}

\subsection{An odd residue disc}
In this subsection \(e=2\), so \(v_l(11)=20\).
Write \(f(X)=X^{11}+b^2\).
\begin{lemma}\label{lem:singulardisc}
Suppose \(P_0=(a,Y_0)\), \(v_l(Y_0)=1\), and \(\bar a=-1\).
If \(P=(a+t,Y(t))\) is on the branch through \(P_0\) and
\(d=v_l(t)\ge28\), then
\[
 [P-P_0]\in H_{d-1}\setminus H_d,
 \qquad
 \log[P-P_0]=\frac{t}{2Y_0}(1,a,a^2,a^3,a^4)
                  \pmod{\varpi^d}.
\]
\end{lemma}
\begin{proof}
Expand the five differentials in the parameter \(t\), on the analytic
branch through \(P_0\). The coefficient of \(t^h\) has valuation at
least \(-2h-1\): use
\(Y(t)^{-1}=Y_0^{-1}(1+(f(a+t)-f(a))/Y_0^2)^{-1/2}\), whose binomial
coefficients are 11-integral. Upon integration, a term with \(h\ge1\)
has valuation at least
\[
 (h+1)d-2h-1-20v_{11}(h+1)>d-1.
\]
Thus the five integrals lie in \(\varpi^{d-1}\OO_l\), and their
leading terms are as asserted. Lemma~\ref{lem:log} allows their
exponential in the formal Picard group. This analytic map and the Abel
map have identical differentials and agree at \(P_0\). Hence they
agree on the disc. The first leading coordinate is nonzero, proving the exact
depth. Note that \(P_0\) need not reduce to a smooth point of the model.
\end{proof}

Write \(H_m^-\) for the anti-invariant subgroup under \(\tau\).
For these subgroups one further has
\begin{equation}\label{eq:H2}
 H_2^- = H_3^-,\qquad \mu_v(H_2^-)=0.
\end{equation}
Indeed, \(\tau\) acts trivially on \(H_2/H_3\), because it fixes the
residue field and sends \(\varpi\) to \(-\varpi\). On an
anti-invariant point it acts by \(-1\). Since 2 is invertible, such a
point is in \(H_3\). The unique formal \(\lambda\)-division in
Lemma~\ref{lem:log} is again anti-invariant. In particular, the local
Kummer class of a point in \(H_1^-\) depends only on its leading
tangent vector in \(H_1/H_2\).

\subsection{The odd rational-image locus}
Suppose \(n\ne n_0\), with \(n\equiv n_0\bmod4356\), and write
\(k_0=v_{11}(n-n_0)\ge2\). Over the completion of \(K\), normalize
\(v_K(u)=1\). The integral minimal elliptic model is
\[
 x'=(x-4)/u^2,\quad y'=(y+1/2)/u^3,\qquad
 (y')^2=(x')^3+\frac{11}{u^2}(x')^2+
           \frac{33}{u^4}x'+\frac{121}{4u^6}.
\]
Its discriminant has valuation three. The formal parameter of \(2G\)
is \(-x'/y'=4u\). The elliptic formal logarithm is an isometry on
this formal neighbourhood, since \(1>v_K(11)/10=1/2\). Thus
\[
 v_K(t'((n-n_0)G))=1+5k_0.
\]
For each odd \(n_0\), put \(q_0=n_0G-3s\). Exact calculations give
\[
 v_K(x'(q_0))=v_K(y'(q_0))=1,\qquad v_K(\delta_{n_0})=2.
\]
The first term in the formal translation of \(x'\) is a unit times
\(2y'(q_0)t'\). Higher terms have larger valuation. Thus
\[
 v_K(x(nG-3s)-x(q_0))=4+5k_0,
\]
and \eqref{eq:cusp} gives
\begin{equation}\label{eq:odddepth}
 v_K(w-1)=5k_0-3,\qquad
 v_{F_v}(X-a)=10k_0-6,\qquad
 v_l(X-a)=20k_0-12.
\end{equation}
Here \(\bar w=1\), and the term \(11(w-1)\) uniquely determines the
valuation of \(w^{11}-1\).
One also checks \(\bar a=-1\) and \(v_l(Y_0)=1\) at the ramified
places. These are the inputs to Lemma~\ref{lem:singulardisc}.

It remains to select a single global square root whose signs agree
with \(P_0\) at both places. Let
\[
 R=((n-n_0)/2)G,\qquad Q=2R,\qquad
 h_\theta(R)=\frac{(x(R)-\theta)^2-(3\theta^2-2\theta-7)}{2y(R)+1}.
\]
Then \(h_\theta(R)^2=x(Q)-\theta\). If
\(m=(y(Q)-y(q_0))/(x(Q)-x(q_0))\), set
\[
 \rho(R)=
 -\frac{m(\theta-x(q_0))+y(q_0)+1/2}
          {(x(q_0)-\theta)h_\theta(R)}.
\]
The exact function-field identities are
\[
 \rho(R)^2=\frac{x(q_0+2R)-\theta}{x(q_0)-\theta},
 \qquad \rho(O)=1.
\]
Choose \(v=v_0\rho(R)\), where \(v_0\) is the known square root in
\eqref{eq:genus5}. This gives the required global choice.
Since \(4\mid n-n_0\), the point \(R\) is formal. Its original elliptic
parameter has \(F_v\)-valuation at least 20, while
\(v_{F_v}(x(q_0)-\theta)\le6\).
The expansion of \(\rho\) at \(O\) therefore has \(\rho(R)-1\) of
positive valuation (in fact at least 14). It chooses the branch through
\(P_0\) simultaneously. We verify the square identity in the elliptic function field.
The value at \(O\) and all the displayed valuations are checked exactly.

Let \(D=[P-P_0]\in J^\eta(F)\). By
Lemma~\ref{lem:singulardisc}, its depth at each ramified place is
\(20k_0-13\). After
\begin{equation}\label{eq:divisions}
 r=10k_0-7\equiv3\pmod{10}
\end{equation}
formal divisions by \(\lambda\), its depth is one. Up to a nonzero
scalar at each place, its leading tangent vector is
\begin{equation}\label{eq:direction}
 \mathbf d=
 \left(1,-\frac1{2^3},\frac1{3^3},-\frac1{4^3},\frac1{5^3}\right)
 =(1,4,9,6,3)\in\F_{11}^5.
\end{equation}
This follows from the diagonal action in Lemma~\ref{lem:log} and
\(\bar a=-1\). The exponent three is essential: the odd elliptic
point has \(v_K(y')=1\), unlike the even cases.

\subsection{The final odd directions}
Here is an explicit computation of the Kummer image of
\eqref{eq:direction}. Choose smooth points with residual
\(X\)-coordinates \(r_i=0,2,3,4,8\) and square roots \(y_i\in k\).
At zero take \(y_i=-b\). The tangent matrix of the five Abel maps is
\[
 M_{ij}=\frac{r_i^{j-1}}{2\bar y_i}.
\]
It is invertible. Solve \(\mathbf uM=\mathbf d\) over \(\F_{11}\)
and lift the \(u_i\) to integers. Let
\[
 P_i^\pm=(r_i\pm\varpi u_i,y_i^\pm),
 \qquad y_i^\pm\equiv y_i\pmod\varpi.
\]
The anti-invariant formal divisor
\(6\sum_i(P_i^+-P_i^-)\) has leading tangent vector \(\mathbf d\),
because \(12\equiv1\bmod11\). Its Kummer image is therefore
\begin{equation}\label{eq:muformula}
 d_v=6\sum_i\bigl(\mu_v(y_i^+-b)-\mu_v(y_i^--b)\bigr).
\end{equation}
By \eqref{eq:H2}, this is the image of every anti-invariant formal
point with that leading tangent vector. The calculations check
\(d_v\ne0\) and membership in the full local image.

Let \(\mathcal S_R^+\) be the restriction of \(\mathcal S^+\) to the
two ramified places, and let
\[
 \mathcal D=\langle(d_{v_1},0),(0,d_{v_2})\rangle.
\]
The exact finite-field certificates give
\begin{center}\begin{tabular}{rcccc}\toprule
\(n_0\)&\(\dim\mathcal S^+\)&localization kernel&\(\dim\mathcal D\)&
\(\dim(\mathcal S_R^+\cap\mathcal D)\)\\\midrule
\(-1\)&1&0&2&0\\
\(1\)&1&0&2&0\\
\(3\)&1&0&2&0\\\bottomrule
\end{tabular}\end{center}
We retain the global and local matrices, as well as the exact
global representatives. Separate finite-field linear algebra
reconstructs the Selmer upper bounds and verifies these intersections.

\begin{proposition}\label{prop:odd}
If \(n_0\in\{-1,1,3\}\), \(n\equiv n_0\bmod4356\), and
\eqref{eq:cusp} holds, then \(n=n_0\).
\end{proposition}
\begin{proof}
Suppose otherwise and form \(D\) as above. Its local Kummer images at
the two ramified places vanish, since its formal depth is at least
three. By \eqref{eq:selmer}, its global Kummer class is zero. Hence
there is a global \(\lambda\)-division.
At each ramified place the division is unique, because the twisted
\(\lambda\)-torsion is zero. It is therefore exactly the formal
division of Lemma~\ref{lem:log}.
Repeat this argument while the depth is at least three. After the
\(r\) divisions in \eqref{eq:divisions}, one obtains a global point
whose pair of local Kummer images is in \(\mathcal S_R^+\).
Equations \eqref{eq:direction}--\eqref{eq:muformula} show that this
pair is a nonzero element of \(\mathcal D\), contradicting the
computed zero intersection.
\end{proof}


\subsection{The even classes}
We now let \(n_0\in\{-2,0,2,4\}\). The curve \eqref{eq:unit} is
already defined over \(F\), and \(J[\lambda](F)=\langle T\rangle\)
has order 11. At the three primes above 11 we have \(e=1\),
\(\bar a=8\), and \(\overline{Y(P_0)}\ne0\). Thus
\begin{equation}\label{eq:even-division}
 \lambda:H_m\xrightarrow{\sim}H_{m+1}\qquad(m\ge1)
\end{equation}
by Lemma~\ref{lem:log}.
At each such place write \(k=F_v\). We use the Abel map
\(i:C\to J\), \(i(P)=[P-P_0]\).

Let \(\mathcal S\) be the localization of the \(\lambda\)-Selmer
group at the three primes over 11. Define
\[
 U=\prod_{v\mid11}\mu_v(H_1),\quad
 T_{\mathrm{loc}}=\prod_{v\mid11}\langle\mu_v(T)\rangle,\quad
 T_{\mathrm{glob}}=\langle(\mu_v(T))_v\rangle.
\]
The exact Selmer computations give
\[
\begin{array}{c|rrrr}
 n_0&-2&0&2&4\\\hline
 \dim\Sel^\lambda(J/F)&5&5&4&4\\
 \dim\ker(\mathrm{loc}_{11})&0&0&0&0\\
 \dim U&3&3&3&3\\
 \dim T_{\mathrm{loc}}&3&3&3&3
\end{array}
\]
and
\begin{equation}\label{eq:intersection}
 U\cap T_{\mathrm{loc}}=0,\qquad
 \mathcal S\cap(U+T_{\mathrm{loc}})=T_{\mathrm{glob}}.
\end{equation}
Every local torsion image is nonzero. The class-group vanishing
\eqref{eq:classes} also applies here. The split descent algebra is
\(F\times F\), whose global \(S\)-class space has dimension 38.
The calculation certifies 38 candidate classes by exact ideal
membership and independent joint local images. It then imposes the
norm and local conditions, including the condition at 2. This
certifies the global space without a GRH assumption. The identities
\eqref{eq:intersection} are verified from the saved localization
matrices by exact linear algebra.

\begin{proposition}\label{prop:even}
If \(n_0\in\{-2,0,2,4\}\), \(n\equiv n_0\bmod4356\), and
\eqref{eq:cusp} holds, then \(n=n_0\).
\end{proposition}

We first describe the local formal image used in
\eqref{eq:intersection}. Here \(\pi=1-\zeta_{11}\) is a uniformizer
and \(11/\pi^{10}\equiv-1\pmod\pi\).
To compute \(\mu_v(H_1)\), choose smooth sections with distinct
residual \(X\)-coordinates
\(0,2,3,4,8\). At \(X=0\) choose \(Y=-b\).
The sum of the five Abel maps is a system of formal coordinates, since
its differential is a Vandermonde matrix with unit determinant.
For a displacement \(X\mapsto X+\pi u\), direct expansion gives
\[
 \frac{Y_{\rm new}-b}{Y_{\rm old}-b}
 \equiv
 \begin{cases}
 1+\pi^{11}\bar u/(4\bar b^2)&\bar X=0,\\
 1&\bar X\in\{2,3,4,8\}
 \end{cases}\pmod{\pi^{12}}.
\]
Indeed, at a nonzero residue the two possible terms of valuation 11
in \((X+\pi u)^{11}-X^{11}\) cancel:
\( (11/\pi^{10})\bar X^{10}\bar u+\bar u^{11}=0\).
All elements of \(1+\pi^{12}\OO_k\) are eleventh powers, by logarithm
and exponential on \(1+\pi^2\OO_k\). The first line generates the
nonzero unramified Kummer line. Its nonvanishing is also verified by
the local Kummer maps in the computation. Thus \(\dim\mu_v(H_1)=1\).

Let
\[
 \ell_0(X)=\frac{(X-2)(X-3)(X-4)(X-8)}{(-2)(-3)(-4)(-8)}
 =1+2X+2X^2+X^3+9X^4\in\F_{11}[X].
\]
This is the linear functional extracting the displacement at the section
with residual coordinate zero, expressed in the differential basis.
For a point \(P\) near \(P_0\), suppose \(v_k(X(P)-X(P_0))=d\ge2\).
Then \(i(P)\in H_d\setminus H_{d+1}\). Put
\(Q=\lambda^{-(d-1)}i(P)\in H_1\), using the unique formal divisions
in \eqref{eq:even-division}. Up to a nonzero scalar, its Kummer image is
\begin{equation}\label{eq:leading}
 h_{d-1}(8),\qquad
 h_r(x)=\sum_{j=1}^5\ell_{0,j-1}\,x^{j-1}j^{-r},
\end{equation}
where \(\ell_{0,j-1}\) is the coefficient of \(X^{j-1}\) in \(\ell_0\).
This follows by applying the diagonal leading action of successive
formal divisions to the Abel tangent vector
\((1,8,8^2,8^3,8^4)/(2\overline{Y(P_0)})\).
All nonlinear terms disappear on the associated graded groups.
In particular,
\begin{equation}\label{eq:five}
 r\equiv1\pmod{10}\quad\Longrightarrow\quad h_r(8)=h_1(8)=5\ne0.
\end{equation}

Normalize \(v_K(u)=1\) as above, and use the same integral minimal model.
Its discriminant has valuation three. The point \(2G\) is formal, with
parameter \(-x'/y'=4u\). If \(n\ne n_0\), let
\(k_0=v_{11}(n-n_0)\ge2\). Since \(n-n_0\) is even, the elliptic
formal logarithm, which is an isometry here, gives
\[
 v_K(t((n-n_0)G))=1+5k_0.
\]
For \(q_0=n_0G-3s\), the exact computations give
\[
 v_K(x(q_0)-4)=2,\quad v_K(y(q_0)+1/2)=3,
 \quad v_K(\delta_{n_0})=2.
\]
Thus \(q_0\) is a smooth point on the special fibre of the minimal
model, with unit \(y'\), and translation by the formal point gives
\[
 v_K(x(nG-3s)-x(q_0))=3+5k_0.
\]
Equation \eqref{eq:cusp} now yields
\[
 v_K(w^{11}-1)=1+5k_0,
 \qquad v_K(w-1)=5k_0-4.
\]
For the second equality, first reduce modulo \(u\) to get \(\bar w=1\).
In the binomial expansion the term \(11(w-1)\) then has strictly
smaller valuation than all other nonzero terms, since \(v_K(11)=5\).
At every \(v\mid11\) of \(F\) we obtain
\begin{equation}\label{eq:d}
 d=v_k(X(P)-X(P_0))=10k_0-8\ge12,
 \qquad d-1\equiv1\pmod{10}.
\end{equation}

We choose the global square root so that \(P\) is near \(P_0\) at
all three places. Put
\[
 R=(n/2)G+4s,\qquad R_0=(n_0/2)G+4s,
\]
so \(2R=nG-3s\). Then
\begin{equation}\label{eq:root}
 v(R)=\frac{(x(R)-\theta)^2-(3\theta^2-2\theta-7)}{2y(R)+1}
\end{equation}
satisfies \(v(R)^2=x(2R)-\theta\). At \(R_0\) this equals the chosen
known square root exactly. Since \(4\mid M\), the difference
\(R-R_0=((n-n_0)/2)G\) is formal. Its parameter has valuation at least
22 over each \(F_v\). All \(R_0\) have integral minimal coordinates.
The computed valuations are
\[
 4\le v_k(x(R_0)-\theta)\le6,\quad
 v_k(2y(R_0)+1)\le8,\quad v_k(v(R_0))=2.
\]
Translation in the minimal coordinates changes the original \(x\)
and \(y\) coordinates by quantities of valuations at least 26 and 28.
Substituting in \eqref{eq:root} shows
\(v(R)/v(R_0)\equiv1\pmod{\pi^{20}}\).
This fixes the required signs simultaneously. Since \(2Y(P_0)\) is
a unit, \(P\) is in its smooth residue disc, and \eqref{eq:d} applies.

\begin{lemma}\label{lem:global}
Under the Selmer injectivity and intersection conditions in
\eqref{eq:intersection}, no global point \(D\in J(F)\) can satisfy
\(D\in\prod_{v\mid11}H_{d,v}\), with \(d\ge2\), and have a nonzero
Kummer image after its \(d-1\) successive formal \(\lambda\)-divisions.
\end{lemma}
\begin{proof}
Write \(Q_j\) for the formal quotient \(\lambda^{-j}D\), for
\(0\le j\le d-1\). Initially \(Q_0=D\) is global, and its local
Kummer image is zero because \(H_d\subset\lambda H_1\).
Injectivity of Selmer localization implies that its global Kummer
class is zero. Hence a global \(R\) with \(\lambda R=Q_0\) exists.
Locally, \(R=Q_1+t_vT\), for scalars \(t_v\in\F_{11}\). Its local
Kummer image lies both in \(\mathcal S\) and in \(U+T_{\rm loc}\).
By \eqref{eq:intersection}, the formal contribution is zero and all
\(t_v\) coincide. Subtracting this global multiple of \(T\) shows
that \(Q_1\) is global and has zero global Kummer class, the latter
again by localization injectivity. Repeating the argument shows that
the local Kummer image of \(Q_{d-1}\) is zero, as asserted.
\end{proof}

\begin{proof}[Proof of Proposition~\ref{prop:even}]
If \(n\ne n_0\), the construction above gives \(D=i(P)\in J(F)\)
with depth \(d=10k_0-8\) at every place above 11.
Equations \eqref{eq:leading}--\eqref{eq:five} show that its final formal
quotient has nonzero Kummer image at every such place. This contradicts
Lemma~\ref{lem:global}.
\end{proof}

This is the division-image argument of Selmer-group Chabauty,
applied to the isogeny \(\lambda\). See
\cite[Theorem~2.1 and Remark~2.2]{Stoll}.


\begin{proof}[Proof of Theorem~\ref{thm:main}]
The six-prime sieve puts the image of every rational point into one
of the seven states \eqref{eq:sieve}. The cusp identity gives
\eqref{eq:cusp}. Propositions~\ref{prop:odd} and \ref{prop:even}
then give \(n=n_0\). The seven base points are listed in
Table~\ref{tab:121}, and each has its unique rational CM lift.
There are no rational cusps, since their images would be rational
cusps on \(E\), whose cusp divisor has degree five and is irreducible.
This proves the assertion.
\end{proof}

\section{Computational data}\label{sec:computations}

We use Magma V2.30-1 \cite{Magma} for the algebraic computations
and Python for the finite sieve. The ancillary data contain the
exact inputs, completed logs and the finite-field matrices used
in the proof. We describe the necessary calculations and the
supplementary checks separately.

\subsection{The sieve and the auxiliary curves}
We construct the six-prime tables directly from the normalized
functions \(f_m\). The program \texttt{six\_prime\_sieve.py}
enumerates the possible classes and congruences. It retains every
cusp residue, so no additional criterion at a cusp is needed.
The exact computation of the known CM classes gives the seven
states in \eqref{eq:sieve}.

The files \texttt{verify\_descent.m} and
\texttt{verify\_genus5.m} check the normalized function identity,
the number fields and the seven genus-five equations. The
valuation and square-root identities used in the local argument
are checked by \texttt{verify\_even\_geometry.m},
\texttt{verify\_odd\_geometry.m} and
\texttt{verify\_odd\_depth.m}.

\subsection{The Selmer calculations}
For each odd class, we give fifteen exact elements of \(B\)
and their joint local images. We check that the valuations
outside \(S\) are divisible by 11. Their independent local
images, together with the proved global dimension, certify the
global basis. We generate each of the three local Kummer images
using actual local divisors until its dimension is five. The
resulting matrices determine the Selmer upper bound and verify
\eqref{eq:selmer} and the zero intersections with the two formal
direction lines.

For each even class, the data give the Selmer localization, the
three formal lines and the three local torsion lines. Exact
linear algebra verifies \eqref{eq:intersection}. The global
basis certificate and the local conditions at 2 and 11 are
included in the reproduction scripts.

GRH is used to find candidate generators, but each candidate is
subsequently checked by ideal factorization and local
independence. The required dimensions follow from
\eqref{eq:classes}. Thus the global bases and the resulting
Selmer bounds are certified without GRH.

The odd calculation is recorded at precisions 80 and 160. We
accept a local Kummer evaluation only when its relative
precision is at least \(2v(11)+2\), which exceeds the
eleventh-power detection threshold. The finite-field matrix
audits do not depend on a further local approximation. The
formal-group lemmas in Section~\ref{sec:local} give the
mathematical argument relating these data to rational points.

\subsection{Supplementary checks and reproduction}
The cases \(n=n_0+4356\) and \(n=n_0+11\cdot4356\) give
additional checks of the odd depth and direction formulas at
both ramified places. They are checked by
\texttt{verify\_local\_depth\_examples.m}. These examples are
not used in place of the argument for arbitrary
\(v_{11}(n-n_0)\) in Section~\ref{sec:local}.

The command \texttt{audit\_saved\_certificates.py} repeats the
finite sieve, compares its surviving classes with the exact CM
classes and verifies all saved odd and even matrix identities.
It does not repeat the number-field calculations or generate
local points. The ancillary instructions give separate commands
for the full descents and the exact geometry checks, together
with a manifest of SHA-256 hashes.

\begin{thebibliography}{9}

\bibitem{Magma}
W. Bosma, J. Cannon and C. Playoust,
\emph{The Magma algebra system. I. The user language},
J. Symbolic Comput. \textbf{24} (1997), 235--265.

\bibitem{Od}
A. M. Odlyzko,
\emph{Table 4: Unconditional bounds for discriminants},
November 29, 1976,
\url{https://www-users.cse.umn.edu/~odlyzko/unpublished/discr.bound.table4}.

\bibitem{OdDescription}
A. M. Odlyzko,
\emph{Discriminant bounds: description of tables},
November 29, 1976,
\url{https://www-users.cse.umn.edu/~odlyzko/unpublished/discr.bound.tables.txt}.

\bibitem{RSZB}
J. Rouse, A. V. Sutherland and D. Zureick-Brown,
\emph{\(\ell\)-adic images of Galois for elliptic curves over
\(\Q\)}, with an appendix with J. Voight,
Forum Math. Sigma \textbf{10} (2022), e62,
\href{https://arxiv.org/abs/2106.11141}{arXiv:2106.11141}.

\bibitem{ST}
R. Schoof and N. Tzanakis,
\href{https://www.mat.uniroma2.it/~schoof/AA-schoof-tzanakis.pdf}
{\emph{Integral points of a modular curve of level 11}},
Acta Arith. \textbf{152} (2012), 39--49.

\bibitem{Stoll}
M. Stoll,
\href{https://www.mathe2.uni-bayreuth.de/stoll/papers/ratpts-selmer-2016-12-07.pdf}
{\emph{Chabauty without the Mordell--Weil group}}, preprint (2017).
\end{thebibliography}

\end{document}
